\section{The missing split local chart from the actual crystalline component}
\label{sec:cubic-split-crystalline-chart}

The local transport of AC-G49-inert did not cover a split
coefficient prime with nontrivial local cubic character.  Here the
local multiplier is trivial, so the actual local problem is a framed
$\mathrm{GL}_3$ problem.  We construct its selected component directly.
The published input is Kisin, \emph{The structure of potentially
semi-stable deformation rings}, ICM 2010, \S1.1.4, equation (1.1.6),
and Theorem 1.2.2, pp.~5--7 of the author version
\url{https://people.math.harvard.edu/~kisin/dvifiles/licm.pdf}.
The latter theorem treats arbitrary residual framed representations
and all finite local coefficient algebras, including nonreduced ones.
It is an established local deformation theorem used here separately
from A; it is not inferred from the definition
\texttt{KisinCrFull} in the source archive.

Let $\ell\ne2,3,269$ split in $K=\mathbb Q(\omega)$, put
$E=\operatorname{Frac}(O)$, and retain the coefficient ring and
the old local and \texttt{GD} conditions.  On the chosen rational
local decomposition group put $\xi=\chi^j$, $j\in\{1,2\}$.
It is an unramified character with values in $O^\times$; it need
not be trivial.  The actual integral local basepoint is
\[
 r_*(g)=\operatorname{diag}
           (\xi(g)\gamma(g),\ 1,\ \xi(g)\gamma(g)^{-1}),
 \qquad \gamma=\psi\kappa_\ell^e,\quad e\ge1.                 \tag{G54.1}
\]
Write $U$ for the unrestricted framed universal local $O$-algebra
of its residual representation, and $x_*:U\to O$ for this
basepoint.  Its existence is the arbitrary-residue framed construction
already used in AC-G49.  Let $\mathfrak a=\ker x_*$ and use the
weight torus $T(t)=\operatorname{diag}(t,1,t^{-1})$.

\begin{theorem}[AC-G54: actual split cubic local component and graded chart]
\label{AC-G54}
There is a prime ideal $P_*\subset\mathfrak a$ and an
$O$-algebra homomorphism
\[
       \Phi:U\longrightarrow E[[X_1,\ldots,X_{12}]]           \tag{G54.2}
\]
with the following properties.
\begin{enumerate}
\item Its kernel is $P_*$ and its constant coefficient is $x_*$.
There are weights $w_i\in\{-2,-1,0,1,2\}$, of multiplicities
\[
              (1,2,3,4,2),                                  \tag{G54.3}
\]
such that every actual universal matrix entry has weight
$\operatorname{wt}(a,b)=\operatorname{wt}_1(a)-
\operatorname{wt}_1(b)$ under $\Phi$.
Its linear coefficients span the twelve coordinates.
If $x_*(f)=0$ and all linear coefficients of $\Phi(f)$ vanish,
then some $0\ne c\in O$ satisfies
$cf\in\mathfrak a^2+P_*$.
\item The condition of killing $P_*$ is stable under every
strict conjugation over every complete Noetherian local
$O$-test algebra with the original residue field, including
algebras with $\ell$-torsion and nilpotents.
\item Every finite local $E$-algebra point killing $P_*$ is
potentially crystalline of the same Galois and Hodge type as
(G54.1).  In particular this holds for every dual-number point
of that component.  Conversely, every such point specializing
to $x_*$ kills $P_*$.
\end{enumerate}
Thus this is the full actual \texttt{LocGood} input, with the
same dimensions as AC-G49-inert, at the previously uncovered
split coefficient primes.  It is a selected component through
the cubic basepoint, not a scalar identification with the old
universal ring.
\end{theorem}
\begin{proof}
\proofstep{The actual fixed-type quotient.}
Take the Galois and Hodge types of (G54.1).  The Hodge weights
are the three distinct weights of $\kappa_\ell^e,1,
\kappa_\ell^{-e}$.  The finite character $\psi$ makes the
representation potentially crystalline, and the unramified
factor $\xi$ does not change its inertial type.
These types are defined over $E$.
Kisin's Theorem 1.2.2 supplies the $\ell$-torsion-free
fixed-type quotient of the framed ring.  Denote its kernel in
$U$ by $I_{\mathrm{cr}}$ and its generic fibre by $B$.
For every finite local $E$-algebra $C$, an $O$-algebra map
$U\to C$ kills $I_{\mathrm{cr}}$ if and only if its
representation has this potentially crystalline type.
In particular $I_{\mathrm{cr}}\subset\mathfrak a$.

All finite Artin points at $x_*$ used below give continuous
representations.  To see the topology explicitly, choose
finitely many topological $O$-algebra generators of $U$ and
subtract their $x_*$-values.  The resulting presentation is
by power series whose variables have zero value at $x_*$.
Their images in an Artin point at $x_*$ are nilpotent, so
evaluation reduces to finitely many coefficients in each
nilpotent degree and is $\ell$-adically continuous.  Explicitly,
the map factors through $O[[Y]]/(Y)^N$ for one fixed $N$;
each matrix coordinate is a finite sum of the continuous
$O$-coefficient maps multiplied by the fixed images of
monomials of degree less than $N$.
The same argument works after adjoining finitely many formal
conjugation variables and truncating them.

\proofstep{The full local ring is regular at the basepoint.}
The six off-diagonal characters of $\operatorname{ad}r_*$ are
\[
\begin{array}{c|ccccc}
\text{weight}&-2&-1&0&1&2\\ \hline
\text{coefficient}&\gamma^{-2}&
 \xi^{-1}\gamma^{-1}\oplus\xi\gamma^{-1}&
 1^{\oplus3}&\xi\gamma\oplus\xi^{-1}\gamma&\gamma^2.
\end{array}                                                  \tag{G54.4}
\]
The nonzero-weight characters have no invariants, since a
nonzero cyclotomic power has infinite image on every open
inertia subgroup and the remaining factors have finite image.
Thus $H^0(\operatorname{ad}r_*)$ has dimension three and
weight zero.

For potentially crystalline self-extensions of a fixed
type, Kisin's \S1.1.4 identifies the tangent complex with
the two-term crystalline complex and gives
\[
 \dim_E H^1_{\mathrm{pcr}}(\operatorname{ad}r_*)
 =\dim_E D_{\mathrm{dR}}(\operatorname{ad}r_*)/
                 \operatorname{Fil}^0+\dim_E H^0=3+3=6.
                                                               \tag{G54.5}
\]
The same complex applied to each summand in (G54.4) gives
one finite extension dimension for each positive-weight
character, zero for each negative-weight character, and
one for each trivial diagonal character.  Here positive
weight means a positive power of $\kappa_\ell$, as in the
source's weight convention.  The three positive roots
therefore account for the first term in (G54.5).

Framed tangent vectors are cocycles, not cohomology classes.
The coboundaries have dimension $9-3=6$, one in each
off-diagonal matrix direction.  Their sum with the finite
cohomology dimensions gives respectively
\[
 \dim_E Z^1_{\mathrm{pcr}}(\operatorname{ad}r_*)_i
  =1,2,3,4,2\qquad(i=-2,-1,0,1,2).                           \tag{G54.6}
\]
In particular the embedding dimension of the full local
ring $B_{x_*}$ is twelve.

Kisin's theorem says that its irreducible components are
formally smooth of dimension $3^2+3=12$.  A component through
$x_*$ therefore gives a quotient of the local ring with
dimension twelve.  Krull dimension is at most embedding
dimension, so $B_{x_*}$ has both dimensions twelve and is
a regular local ring.  This step uses the tangent calculation;
smoothness of the individual components alone would not
justify regularity at a possible intersection.
Its completion is consequently $E[[T_1,\ldots,T_{12}]]$.
It has a unique minimal prime at $x_*$.

\proofstep{Select the integral prime and retain the saturated chart.}
Compose $U\to B_{x_*}$ with its completion.  Let $P_*$ be
the kernel and $\Phi_0$ this composite, using the displayed
power-series coordinates.  Since the target is a domain,
$P_*$ is prime; also $I_{\mathrm{cr}}\subset P_*\subset
\mathfrak a$.
The elementary completion argument in
\path{PadicHodgeTheory/SmoothCompl.lean},
\texttt{smoothCompl}, applies to this formally smooth
Artin lifting problem.  It gives the full
\texttt{IsCompl} properties: constant term $x_*$,
density after multiplication by a nonzero element of $O$,
and, for every $n$,
\[
 \Phi_0(f)\in(T)^n\quad\Longrightarrow\quad
 \exists\,0\ne c\in O:\ cf\in\mathfrak a^n+P_*.
                                                               \tag{G54.7}
\]
It also gives an element $y\in U$ such that
\[
 x_*(y)\ne0,\qquad yP_*\subset I_{\mathrm{cr}}.              \tag{G54.8}
\]
For completeness, (G54.7) is the usual denominator clearing
in $U[1/\ell]$ at the $E$-valued point.  Localization outside
that point does not change the quotient by
$\mathfrak a^n+P_*$ after $\ell$ is inverted: any inverted
element has nonzero constant term and an inverse given by
a finite nilpotent geometric series in this quotient.
The resulting finite denominator is an element of $O$.
Likewise $P_*/I_{\mathrm{cr}}$ vanishes after localization
at $x_*$; finite generation gives a single $y$ as in
(G54.8).  This also explains why one needs only
$x_*(y)\ne0$, not that it be an integral unit.

It follows at once from (G54.8) that every Artin fixed-type
point at $x_*$ kills $P_*$: its value on $y$ has nonzero
constant term and hence is a unit.  Conversely killing
$P_*$ implies killing $I_{\mathrm{cr}}$.  This proves
property 3, with the stated specialization condition in
the converse.

\proofstep{Strict conjugation over arbitrary integral test algebras.}
Put $R=U/P_*$.  It is an $O$-torsion-free domain, and
$\Phi_0$ induces an injection $R\hookrightarrow E[[T]]$.
Adjoin nine formal variables $Z_{ab}$ to $R$.
The matrix $h=I+(Z_{ab})$ is invertible in $R[[Z]]$ and
is strict.  Conjugating the actual universal representation
over $R[[Z]]$ therefore gives, by the unrestricted framed
universal property, a homomorphism
\[
                      a:U\longrightarrow R[[Z]].             \tag{G54.9}
\]
Embed its target coefficientwise into $E[[T,Z]]$.
Modulo any power of the ideal $(T,Z)$, the resulting
representation is conjugate to a fixed-type Artin point
at $x_*$.  The original point kills $I_{\mathrm{cr}}$,
so Kisin's full Artin statement and invariance under
conjugation show that the conjugated point also kills
$I_{\mathrm{cr}}$.
Separatedness of $E[[T,Z]]$ then gives
$a(I_{\mathrm{cr}})=0$ in that ring.

The value of $a(y)$ has constant term $x_*(y)\ne0$,
so is a unit in $E[[T,Z]]$.
Equation (G54.8) now implies $a(P_*)=0$ there.
The coefficientwise embedding of $R[[Z]]$ is injective,
so $a(P_*)=0$ already in $R[[Z]]$.

Let $C$ be any allowed complete local integral test
algebra, let $f:U\to C$ kill $P_*$, and let
$h_C\in\mathrm{GL}_3(C)$ satisfy $h_C-I\in M_3(\mathfrak m_C)$.
The map $R\to C$ and the assignments
$Z_{ab}\mapsto(h_C-I)_{ab}$ define a convergent evaluation
$R[[Z]]\to C$.  Evaluating (G54.9) gives exactly the
conjugated representation, so it kills $P_*$.
No injection from $C$ to a characteristic-zero algebra
has been assumed.  This proves property 2 even when
$C$ has torsion or nilpotents.

\proofstep{Homogeneous coordinates with the exact weights.}
Choose $t=1+\ell\in O^\times$.  It is congruent to one
and has infinite multiplicative order in characteristic
zero.  Conjugation by $T(t)$ induces an automorphism
$\theta$ of $U$, with
\[
 \theta(r^{\mathrm{univ}}(g)_{ab})
          =t^{\operatorname{wt}(a,b)}r^{\mathrm{univ}}(g)_{ab}.
                                                               \tag{G54.10}
\]
The diagonal basepoint is fixed, and the strict-conjugation
result gives $\theta(P_*)=P_*$, using the inverse
conjugation for the reverse inclusion.
The actual entry differences
$e_{g,a,b}=r^{\mathrm{univ}}(g)_{ab}-x_*(r^{\mathrm{univ}}(g)_{ab})$
are eigenvectors of the same weights.  A subtracted
constant occurs only when $a=b$, whose weight is zero.

Their linear coefficients span the cotangent space:
a derivation annihilating every actual entry gives the
zero framed dual-number deformation, hence is zero by
the unrestricted universal property.  This is exactly
the general \texttt{chartSpan} argument in
\path{PadicHodgeTheory/ChartSpan.lean}; denominators can
be cleared to compare the dual-number maps over $O$.
Select twelve actual entry differences whose classes
are a basis.  They are formal parameters in the regular
completed local ring.  The inverse-function construction
\texttt{chartChange} in
\path{PadicHodgeTheory/ChartChange.lean} changes
$\Phi_0$ to a chart $\Phi$ sending these parameters
exactly to $X_1,\ldots,X_{12}$, preserving (G54.7).

Let $w_i$ be the weight of the selected entry difference.
Equation (G54.10) and uniqueness in the completion show
that $\theta$ acts there by $X_i\mapsto t^{w_i}X_i$.
If an entry has weight $j$, comparison of the coefficient
of $X^\nu$ gives
\[
   (t^{\sum_iw_i\nu_i}-t^j)[X^\nu]\Phi(r^{\mathrm{univ}}(g)_{ab})=0.
                                                               \tag{G54.11}
\]
Distinct integral powers of $t$ are distinct, so only
monomials of total weight $j$ occur.  This is the full
generic \texttt{chartGraded} argument of
\path{PadicHodgeTheory/ChartGraded.lean}.
The tangent weights in (G54.6) force exactly the
multiplicities (G54.3).  The chosen parameters give
the onto-linear-coordinate condition, and (G54.7)
with $n=2$ gives the asserted cotangent injectivity.
All fields of property 1 are established.
\end{proof}

\begin{corollary}[AC-G54-model: the cubic model and integral extraction at every good coefficient prime]
\label{AC-G54-model}
Retain the old \texttt{GD}, \texttt{CohInp}, full
\texttt{LemmaP} and, for the length comparison,
\texttt{LemmaM'} assumptions of AC-G49--G52.
For every $\ell\ne2,3,269$, the actual cubic construction
has a complete graded model and
\[
                  \operatorname{length}_O\mathcal S_j
                                \le\mathcal M_j.\mathrm{len}.
                                                               \tag{G54.12}
\]
At the previously uncovered split primes its local ideal
is the prime $P_*$ of AC-G54.
\end{corollary}
\begin{proof}
For inert primes, or split primes with trivial local
$\chi$, use the transported local ideal of AC-G49-inert.
For the remaining split primes use AC-G54.
In both cases the unrestricted local representation
has the actual cubic diagonal basepoint, the local
ideal is strictly conjugation-stable over every
\texttt{IsCO} test algebra, and its chart has the same
twelve weights.

These are the entire local inputs to AC-G49's global
framed construction, AC-G50's local homogeneous
presentation and relative equation count, and AC-G51's
graded model and integral derivative lift.  The global
residual inertia kernels, order-prime-to-$\ell$ argument,
cyclotomic witness and eight normalization equations
are unchanged.  The local presentations in those proofs
use the arbitrary-seed universal property and strict
weight-torus conjugation, both supplied here.
Thus their complete constructions apply to this
actual $P_*$, without an identification with the old
local residual representation.

Finally AC-G52 only requires the constructed full model,
strict local conjugation stability, the old marked unit,
and an inertia element on which the cubic character is
trivial and $\gamma^2-1$ is a unit.  Since
$\ell\ne3,269$, $\chi$ is unramified at $\ell$, so the
same old \texttt{GD} element supplies the last condition.
Every step of its integral cocycle and square-zero
construction therefore applies, giving (G54.12).
\end{proof}

\paragraph{Remaining scope.}
This removes the split nontrivial-character local-chart
gate.  It does not prove a finite bound for the right
side of (G54.12), $M2$, Type I automorphy, or mixed
tangent vanishing.  The ordinary-inertia description
of \emph{all} points of the newly selected component
also requires its own verification before importing
every pointwise Type II argument from the old
$I_Y$-based construction.  Potential crystallinity
of its actual dual-number points, as stated above,
is already proved.  The excluded coefficient primes
remain outside the corollary.

\paragraph{Geometric dictionary.}
The new object is an explicit irreducible component
of a framed local deformation scheme, with a regular
formal germ and a weight action through its actual
matrix entries.  Its finite inertial type is a
finite local system in the audited SGA~1, Expos\'e~V,
\S\S4 and~7 framework.  Its crystalline equations
and formal smoothness come from the precisely cited
Kisin theorem, and its integral conjugation stability
from the argument (G54.8)--(G54.9).  No unexamined
archival scan or general existence of a motivic
realization is used to supply these properties.
